\documentclass[11pt, a4paper]{article}

% ─── Encodage et langue ───
\usepackage[utf8]{inputenc}
\usepackage[T1]{fontenc}
\usepackage[french]{babel}

% ─── Maths ───
\usepackage{amsmath, amssymb, amsthm}

% ─── Mise en page ───
\usepackage[margin=2.5cm]{geometry}
\usepackage{enumitem}
\usepackage{booktabs}
\usepackage{multirow}
\usepackage{float}
\usepackage{xcolor}
\usepackage{colortbl}
\usepackage{tikz}
\usetikzlibrary{arrows.meta, positioning, fit, backgrounds, calc, decorations.pathreplacing}
\usepackage{pifont}
\usepackage{hyperref}
\hypersetup{colorlinks=true, linkcolor=blue!60!black, urlcolor=blue!70!black}
\usepackage{fancyhdr}

% ─── En-tête / pied de page ───
\pagestyle{fancy}
\fancyhf{}
\fancyhead[L]{\small\textsc{Rapport de synthèse}}
\fancyhead[R]{\small\textsc{Mars 2026}}
\fancyfoot[C]{\thepage}

% ─── Commandes ───
\newcommand{\R}{\mathbb{R}}
\newcommand{\Ltwo}{L^2(\mathcal{T})}
\newcommand{\Dp}{D_p}
\newcommand{\Dw}{D_w}
\newcommand{\DK}{D_K}

% ─── Couleurs ───
\definecolor{func}{HTML}{2166AC}
\definecolor{vect}{HTML}{B2182B}
\definecolor{mix}{HTML}{4DAF4A}
\definecolor{gris}{HTML}{666666}
\definecolor{emphbox}{HTML}{FFF3CD}
\definecolor{emphborder}{HTML}{856404}
\definecolor{resultbox}{HTML}{D4EDDA}
\definecolor{resultborder}{HTML}{155724}

% ─── Boîtes d'emphase ───
\usepackage{framed}

% keybox : fond jaune, cadre ambre
\newenvironment{keybox}[1][]{%
  \def\FrameCommand{\fboxsep=8pt \fcolorbox{emphborder}{emphbox}}%
  \MakeFramed{\advance\hsize-\width \FrameRestore}%
  \noindent\textbf{#1}\par\smallskip\noindent\ignorespaces
}{\endMakeFramed}

% resultbox : fond vert, cadre vert foncé
\newenvironment{resultbox}[1][]{%
  \def\FrameCommand{\fboxsep=8pt \fcolorbox{resultborder}{resultbox}}%
  \MakeFramed{\advance\hsize-\width \FrameRestore}%
  \noindent\textbf{#1}\par\smallskip\noindent\ignorespaces
}{\endMakeFramed}

% ============================================================================
\title{%
  \vspace{-1cm}
  \textbf{Classification non supervisée} \\[0.3em]
  \textbf{de données mixtes fonctionnelles + vectorielles} \\[0.5em]
  \Large Synthèse des travaux — Architecture des distances, \\
  résultats multi-dataset et perspectives
}
\author{Pierre Chambet}
\date{Mars 2026}

\begin{document}
\maketitle
\thispagestyle{fancy}

\begin{abstract}
Ce rapport synthétise l'ensemble des travaux réalisés sur la classification
non supervisée de données mixtes (fonctionnelles + vectorielles). On y
présente le problème, l'architecture des distances construite à deux niveaux,
les résultats obtenus sur quatre jeux de données labellisés, et les
conclusions méthodologiques qui en découlent — notamment la découverte que
les dérivées sont systématiquement informatives lorsqu'on optimise par ARI
($\alpha_{\text{ARI-optimal}} \geq 0.6$ sur 4/4 datasets) et l'identification du critère de sélection des paramètres
comme verrou principal.
\end{abstract}

\tableofcontents
\newpage

% ============================================================================
\section{Le problème : mesurer la similarité entre objets mixtes}
% ============================================================================

Chaque individu $i = 1, \ldots, n$ possède deux types d'information
qui vivent dans des espaces mathématiques \textbf{différents} :

\[
  \text{Individu } i : \quad
  \underbrace{\hat{X}_i(t) \in \Ltwo}_{\text{courbe fonctionnelle}}
  \quad + \quad
  \underbrace{Z_i = (z_{i1}, \ldots, z_{ip}) \in \R^p}_{\text{vecteur classique}}.
\]

\begin{figure}[H]
\centering
\begin{tikzpicture}[
    box/.style={draw, rounded corners=6pt, minimum width=4.5cm, minimum height=3cm,
                align=center, font=\small},
  ]
  \node[box, fill=func!8] (ind1) {
    \textbf{Individu $i$} \\[6pt]
    Courbe $X_i(t)$ \\
    {\footnotesize vit dans $L^2$} \\[6pt]
    $Z_i = (48.4, -123.3, 1.6)$ \\
    {\footnotesize vit dans $\R^3$}
  };
  \node[box, fill=vect!8, right=2.5cm of ind1] (ind2) {
    \textbf{Individu $j$} \\[6pt]
    Courbe $X_j(t)$ \\
    {\footnotesize vit dans $L^2$} \\[6pt]
    $Z_j = (43.1, -79.4, 2.8)$ \\
    {\footnotesize vit dans $\R^3$}
  };
  \draw[{Stealth[length=4mm]}-{Stealth[length=4mm]}, thick, dashed, gris]
    (ind1) -- (ind2) node[midway, above, font=\large\bfseries] {$d(i,j) = \;?$};
\end{tikzpicture}
\caption{Le problème central : construire une distance entre objets
qui vivent dans des espaces différents ($L^2$ et $\R^p$).}
\end{figure}

On ne peut pas empiler une courbe et un vecteur dans un même objet, ni
calculer une distance euclidienne entre eux. Il faut une \textbf{architecture}
de distances qui :

\begin{enumerate}[nosep]
  \item Mesure la dissimilarité \textbf{entre courbes} (fonctionnel).
  \item Mesure la dissimilarité \textbf{entre vecteurs} (classique).
  \item \textbf{Combine} les deux de manière contrôlée, avec des paramètres
    interprétables.
\end{enumerate}

La solution est construite en \textbf{deux niveaux} hiérarchiques, détaillés
dans les sections suivantes.

% ============================================================================
\newpage
\section{Les briques de base (Niveau 0)}
% ============================================================================

Trois distances élémentaires, chacune capturant un aspect différent des données.

% ────────────────────────────────────────────────────────────────────────
\subsection{$D_0$ : distance $L^2$ entre courbes — le \textbf{niveau}}

\[
  D_0(i,j) = \sqrt{\int_{\mathcal{T}} \bigl[\hat{X}_i(t) - \hat{X}_j(t)\bigr]^2 \, dt}
\]

$D_0$ mesure la \textbf{proximité globale des profils} : c'est l'aire entre
les deux courbes (au sens $L^2$). Deux courbes au même niveau mais
légèrement décalées en phase seront proches en $D_0$.

\begin{itemize}[nosep]
  \item[\textcolor{mix}{$\boldsymbol{+}$}] Interprétation directe : ``l'écart vertical moyen''.
  \item[\textcolor{vect}{$\boldsymbol{-}$}] Insensible aux différences de
    \textbf{forme} si les niveaux sont proches.
\end{itemize}

\paragraph{Exemple.}
Vancouver et Montréal ont des formes de courbes similaires (hiver froid, été
chaud) mais Montréal est décalée vers le bas. $D_0$ les sépare bien. En
revanche, Vancouver (courbe plate) et Halifax (forte amplitude) ont des
niveaux proches mais des dynamiques très différentes — $D_0$ ne le voit pas.

% ────────────────────────────────────────────────────────────────────────
\subsection{$D_1$ : distance $L^2$ entre dérivées — la \textbf{forme}}

\[
  D_1(i,j) = \sqrt{\int_{\mathcal{T}} \bigl[\hat{X}_i'(t) - \hat{X}_j'(t)\bigr]^2 \, dt}
\]

$D_1$ mesure la proximité des \textbf{dynamiques}. Si $X_i(t)$ est une
température, alors $X_i'(t)$ est la vitesse de changement.

\begin{itemize}[nosep]
  \item[\textcolor{mix}{$\boldsymbol{+}$}] Capture la \textbf{forme},
    indépendamment du niveau vertical.
  \item[\textcolor{vect}{$\boldsymbol{-}$}] Sensible au bruit (la dérivée
    amplifie les fluctuations). Nécessite un bon lissage en amont.
\end{itemize}

\paragraph{Exemple.}
Vancouver (Pacifique) : hivers doux, étés doux $\Rightarrow$ dérivée
\textbf{plate}. Halifax (Atlantique) : hivers froids, étés chauds
$\Rightarrow$ dérivée \textbf{forte} en mars et octobre. $D_0$ les voit
proches (niveaux similaires). $D_1$ les voit \textbf{très différentes}.

% ────────────────────────────────────────────────────────────────────────
\subsection{$D_s$ : distance euclidienne sur $Z$ — le \textbf{vectoriel}}

\[
  D_s(i,j) = \left\| Z_i^{\text{std}} - Z_j^{\text{std}} \right\|_2
\]

où $Z^{\text{std}}$ est la standardisation variable par variable
(moyenne~0, variance~1). $D_s$ mesure la proximité dans l'espace des
variables classiques et \textbf{ignore totalement les courbes}. Cette
standardisation met toutes les composantes de $Z$ sur une échelle
comparable avant qu'elles soient combinées au fonctionnel dans les
stratégies mixtes.

% ────────────────────────────────────────────────────────────────────────
\subsection{Bilan : aucune ne suffit seule}

\begin{table}[H]
\centering
\begin{tabular}{lll}
  \toprule
  \textbf{Distance} & \textbf{Ce qu'elle voit} & \textbf{Ce qu'elle rate} \\
  \midrule
  $D_0$ & Niveau vertical des courbes & Forme / dynamique \\
  $D_1$ & Forme / vitesse de changement & Niveau vertical \\
  $D_s$ & Position dans $\R^p$ & Les courbes entièrement \\
  \bottomrule
\end{tabular}
\caption{Complémentarité des trois distances de base.}
\end{table}

Aucune des trois ne capture toute l'information $\Rightarrow$ il faut les
\textbf{combiner}.

% ============================================================================
\newpage
\section{$\Dp(\alpha)$ : combiner niveau et forme (Niveau 1)}
% ============================================================================

\subsection{Le problème d'échelle}

$D_0$ et $D_1$ sont sur des échelles très différentes. Sur Canadian Weather :
\[
  D_0 \in [3.9,\; 532] \qquad D_1 \in [0.51,\; 6.6] \qquad \text{(facteur $\times 80$)},
\]
où le minimum de $D_0$ exclut la diagonale (distance à soi-même nulle).

Si on les additionne brut, $D_1$ ne pèse rien. \textbf{Solution :}
normaliser par le maximum :
\[
  \tilde{D}_0 = \frac{D_0}{\max(D_0)} \in [0,1], \qquad
  \tilde{D}_1 = \frac{D_1}{\max(D_1)} \in [0,1].
\]

\subsection{La formule}

\begin{keybox}[Définition de $\Dp(\alpha)$]
\[
  \Dp(\alpha) = \sqrt{(1 - \alpha)\,\tilde{D}_0^{\,2}
  + \alpha\,\tilde{D}_1^{\,2}},
  \qquad \alpha \in [0, 1]
\]
\end{keybox}

La normalisation de $D_0$ et $D_1$ dans $[0,1]$ avant combinaison
garantit qu'aucune des deux composantes (niveau ou forme) ne domine
mécaniquement la distance. C'est une condition importante pour que les
stratégies mixtes du niveau~2 ($\Dw$ et $\DK$) puissent pondérer
proprement le fonctionnel par rapport au vectoriel.

\subsection{Le curseur $\alpha$}

Le paramètre $\alpha$ agit comme un \textbf{curseur} entre le niveau
et la forme :

\begin{figure}[H]
\centering
\begin{tikzpicture}
  \draw[very thick, func!70] (0,0) -- (10,0);
  \foreach \x/\lab in {0/$\alpha{=}0$, 5/$\alpha{=}0.5$, 10/$\alpha{=}1$} {
    \fill[func] (\x, 0) circle (3pt);
    \node[below=6pt, font=\small] at (\x, 0) {\lab};
  }
  \node[above=8pt, font=\small\bfseries, func] at (0, 0) {Niveau pur};
  \node[above=8pt, font=\small] at (5, 0) {Compromis};
  \node[above=8pt, font=\small\bfseries, func] at (10, 0) {Forme pure};

  \node[below=22pt, font=\footnotesize, gris] at (0, 0) {$\Dp = \tilde{D}_0$};
  \node[below=22pt, font=\footnotesize, gris] at (10, 0) {$\Dp = \tilde{D}_1$};
\end{tikzpicture}
\end{figure}

\subsection{Propriétés}

$\Dp(\alpha)$ est une \textbf{vraie distance} (positivité, symétrie,
inégalité triangulaire par Minkowski sur $\ell^2$ pondéré). C'est important
car PAM nécessite une matrice de distances au sens métrique.

\begin{keybox}[Rapport avec les travaux précédents (Hugo Nguyen, 2023)]
Hugo avait une formule similaire pour $\Dp$, mais combinait uniquement
\textbf{deux distances fonctionnelles} ($D_0$ et $D_1$).
\end{keybox}

% ============================================================================
\newpage
\section{Le mix total : $\Dw$ et $\DK$ (Niveau 2)}
% ============================================================================

On a maintenant $\Dp(\alpha)$ (fonctionnel combiné) et $D_s$ (vectoriel).
Deux stratégies pour les combiner.

\subsection{Architecture complète}

\begin{figure}[H]
\centering
\begin{tikzpicture}[
    base/.style={draw, rounded corners=4pt, minimum width=2.5cm, minimum height=0.9cm,
                 align=center, font=\small\bfseries},
    fleche/.style={-{Stealth[length=3mm]}, thick},
    etiq/.style={font=\footnotesize\itshape, text=gris},
  ]

  \node[base, fill=func!15, text=func] (D0) at (0,0) {$D_0$ \footnotesize niveau};
  \node[base, fill=func!15, text=func] (D1) at (3.5,0) {$D_1$ \footnotesize forme};
  \node[base, fill=vect!15, text=vect] (Ds) at (8.5,0) {$D_s$ \footnotesize vectoriel};

  \node[base, fill=func!25, text=func] (Dp) at (1.75,-1.8) {$\Dp(\alpha)$};
  \node[base, fill=mix!20, text=mix!70!black] (Dw) at (4.0,-3.6) {$\Dw(\alpha, \omega)$};
  \node[base, fill=mix!20, text=mix!70!black] (DK) at (7.5,-3.6) {$\DK(\alpha)$};

  \draw[fleche, func!70] (D0) -- (Dp) node[etiq, midway, left=2pt] {$1{-}\alpha$};
  \draw[fleche, func!70] (D1) -- (Dp) node[etiq, midway, right=2pt] {$\alpha$};
  \draw[fleche, func!70] (Dp) -- (Dw) node[etiq, midway, left=2pt] {$\omega$};
  \draw[fleche, vect!70] (Ds) -- (Dw) node[etiq, midway, right=2pt] {$1{-}\omega$};
  \draw[fleche, func!70] (Dp) -- (DK) node[etiq, midway, left=2pt] {$K_f$};
  \draw[fleche, vect!70] (Ds) -- (DK) node[etiq, midway, right=2pt] {$K_s$};

  \node[font=\footnotesize\scshape, text=gris] at (-1.5,0) {Niv.\ 0};
  \node[font=\footnotesize\scshape, text=gris] at (-1.5,-1.8) {Niv.\ 1};
  \node[font=\footnotesize\scshape, text=gris] at (-1.5,-3.6) {Niv.\ 2};
\end{tikzpicture}
\caption{Architecture à deux niveaux.
\textcolor{func}{Bleu} = fonctionnel,
\textcolor{vect}{rouge} = vectoriel,
\textcolor{mix}{vert} = mixte.}
\end{figure}

Cette architecture peut se lire comme une \textbf{famille d'approches de
type RS-PCA} : la partie fonctionnelle est d'abord réduite dans un
espace de dimension finie (via FPCA et la distance $\Dp(\alpha)$), puis
combinée avec la partie vectorielle $Z$ pour construire soit (i) un
espace euclidien réduit où l'on applique k-means (stratégie A), soit
(ii) des distances mixtes directement dans l'espace d'origine
($\Dw$ et $\DK$, stratégies B et C). La différence principale entre A,
B et C n'est donc pas la réduction elle-même (toujours fondée sur FPCA
et $Z$), mais la manière dont on \textbf{exploite} cette réduction :
par PCA+k-means dans $\R^{K+p}$ pour A, ou par des distances/noyaux
pour B et C.

% ────────────────────────────────────────────────────────────────────────
\subsection{Résumé des trois stratégies de clustering}

\begin{table}[H]
\centering
\renewcommand{\arraystretch}{1.25}
\begin{tabular}{l p{3.2cm} p{3.4cm} p{3.4cm}}
  \toprule
  \textbf{Stratégie} & \textbf{Représentation} &
  \textbf{Méthode de clustering} & \textbf{Rôle de la FPCA / distances} \\
  \midrule
  A (FPCA+Z) &
  Vecteurs $W = [\xi\_{\text{FPCA}} \mid Z]$ standardisés &
  k-means (distance euclidienne) &
  FPCA sur les courbes, $Z$ standardisé, pas d'ACP sur $Z$ \\
  B ($\Dw$) &
  Matrice de distances $\Dw(\alpha,\omega)$ &
  PAM (k-médoïdes) sur distances &
  Combinaison pondérée de $\Dp(\alpha)$ (fonctionnel) et $D_s$ (vectoriel) \\
  C ($\DK$) &
  Matrice de distances $\DK(\alpha)$ &
  PAM (k-médoïdes) sur distances &
  Produit de noyaux issus de $\Dp(\alpha)$ et $D_s$ puis retour à une distance \\
  \bottomrule
\end{tabular}
\caption{Résumé des trois stratégies : type de représentation, méthode de
clustering et usage de la FPCA ou des distances.}
\end{table}

% ────────────────────────────────────────────────────────────────────────
\subsection{Stratégie A : FPCA + $Z$ — concaténation dans $\R^{K+p}$}

La première stratégie mixte correspond à une \textbf{RS-PCA (Reduced-Space
PCA) au sens strict}. La chaîne est la suivante :
\begin{enumerate}[nosep]
  \item lissage des courbes $X_i(t)$ puis FPCA $\Rightarrow$ scores
        fonctionnels $(\xi_{i1}, \ldots, \xi_{iK})$ ;
  \item construction de $Z$ (géographie, spectres agrégés, etc.) et
        standardisation variable par variable ;
  \item concaténation des deux blocs dans un même vecteur ;
  \item standardisation colonne par colonne, puis k-means avec distance
        euclidienne.
\end{enumerate}

Concrètement, avant standardisation finale, on forme pour chaque
individu :
\[
  U_i = (\xi_{i1}, \ldots, \xi_{iK}, Z_i) \in \R^{K+p},
\]
puis on applique k-means sur la matrice $W$ obtenue en concaténant les
scores FPCA et $Z$, tous deux centrés-réduits. Autrement dit, la
distance ``vue'' par la stratégie A est la distance euclidienne dans un
espace réduit où chaque coordonnée (score fonctionnel ou variable de
$Z$) contribue de manière comparable.

Il n'y a \textbf{pas d'ACP supplémentaire sur $Z$} : les variables
vectorielles sont uniquement standardisées. Ce choix est adapté aux cas
étudiés ici, où $Z$ est de petite dimension et doit rester directement
interprétable. La stratégie A sert ainsi de \textbf{baseline mixte}
directe à comparer aux stratégies B et C, qui combinent les sources au
niveau des distances (pondération) ou des noyaux (produit). Dans tout
ce rapport, on se concentre sur cette approche RS-PCA enrichie par des
distances fonctionnelles et mixtes adaptées au clustering et à l'étude
des critères de sélection.

% ────────────────────────────────────────────────────────────────────────
\subsection{Stratégie B : $\Dw(\alpha, \omega)$ — pondération directe}

\begin{keybox}[Définition de $\Dw(\alpha, \omega)$]
\[
  \Dw(\alpha, \omega) = \sqrt{\omega \cdot \Dp(\alpha)^2
  + (1 - \omega) \cdot \tilde{D}_s^{\,2}},
  \qquad \alpha, \omega \in [0, 1]
\]
\end{keybox}

Les deux entrées de cette formule sont donc des \textbf{distances déjà
mise à l'échelle} : $\Dp(\alpha)$, construite à partir de $\tilde{D}_0$
et $\tilde{D}_1$ dans $[0,1]$, et $\tilde{D}_s$, issue de $Z$
standardisé. Le paramètre $\omega$ joue alors le rôle de curseur
explicite entre la contribution fonctionnelle et la contribution
vectorielle, sur une base devenue comparable.

Deux ``boutons'' indépendants :

\begin{figure}[H]
\centering
\begin{tikzpicture}
  % Bouton alpha
  \draw[very thick, func!70] (0, 0.8) -- (5, 0.8);
  \fill[func] (0, 0.8) circle (2.5pt); \fill[func] (5, 0.8) circle (2.5pt);
  \node[left, font=\small\bfseries] at (-0.3, 0.8) {$\alpha$ :};
  \node[above=3pt, font=\footnotesize] at (0, 0.8) {Niveau};
  \node[above=3pt, font=\footnotesize] at (5, 0.8) {Forme};
  \node[right=3pt, font=\footnotesize, gris] at (5.3, 0.8) {(dans le fonctionnel)};

  % Bouton omega
  \draw[very thick, mix!70] (0, 0) -- (5, 0);
  \fill[mix!70!black] (0, 0) circle (2.5pt); \fill[mix!70!black] (5, 0) circle (2.5pt);
  \node[left, font=\small\bfseries] at (-0.3, 0) {$\omega$ :};
  \node[above=3pt, font=\footnotesize] at (0, 0) {Vecteurs $Z$};
  \node[above=3pt, font=\footnotesize] at (5, 0) {Courbes};
  \node[right=3pt, font=\footnotesize, gris] at (5.3, 0) {(entre les 2 sources)};
\end{tikzpicture}
\end{figure}

Cas limites :

\begin{center}
\begin{tabular}{ccl}
  \toprule
  $\boldsymbol{\alpha}$ & $\boldsymbol{\omega}$ & \textbf{Distance effective} \\
  \midrule
  $*$ & $0$ & $\tilde{D}_s$ seul (courbes ignorées) \\
  $0$ & $1$ & $\tilde{D}_0$ seul (niveau des courbes) \\
  $1$ & $1$ & $\tilde{D}_1$ seul (forme des courbes) \\
  $0.5$ & $0.5$ & Compromis global \\
  \bottomrule
\end{tabular}
\end{center}

\textbf{Sélection des paramètres :} grid search 2D sur
$(\alpha, \omega) \in \{0, 0.1, \ldots, 1\}^2$ (121 combinaisons), en
maximisant la \textbf{silhouette moyenne} (critère non supervisé).

% ────────────────────────────────────────────────────────────────────────
\subsection{Stratégie C : $\DK(\alpha)$ — noyaux gaussiens}

Au lieu de pondérer des distances, on passe par des \textbf{similarités}
(noyaux), qu'on combine par produit :

\begin{align}
  K_f(i,j) &= \exp\!\left( -\frac{\Dp(\alpha)_{ij}^2}{2\sigma_f^2} \right),
  & \sigma_f &= \text{médiane}\{\Dp(\alpha)_{ij}\},  \\[4pt]
  K_s(i,j) &= \exp\!\left( -\frac{D_s(i,j)^2}{2\sigma_s^2} \right),
  & \sigma_s &= \text{médiane}\{D_s(i,j)\},  \\[4pt]
  K &= K_f \cdot K_s, & & \text{(produit = ``ET logique flou'')} \\[4pt]
  \DK(i,j) &= \sqrt{K(i,i) + K(j,j) - 2K(i,j)}.
\end{align}

Le choix $\sigma_f = \text{médiane}\{\Dp(\alpha)_{ij}\}$ et
$\sigma_s = \text{médiane}\{D_s(i,j)\}$ joue le rôle de
\textbf{mise à l'échelle automatique} dans chaque noyau : pour chaque
source, des distances typiques correspondent à des similarités de taille
intermédiaire. Le produit $K = K_f \cdot K_s$ signifie ensuite que deux
individus ne sont similaires \textbf{que si les deux sources le
confirment}. Si $K_f \approx 0$ (courbes très différentes), alors
$K \approx 0$ même si $K_s = 1$.

\textbf{Avantage :} pas de paramètre $\omega$ — grid search 1D sur $\alpha$
uniquement (11 valeurs au lieu de 121).

\begin{table}[H]
\centering
\begin{tabular}{lcccc}
  \toprule
  & \textbf{Paramètres} & \textbf{Combinaison} & \textbf{Grid search} \\
  \midrule
  $\Dw$ & $\alpha, \omega$ & Somme pondérée & 2D (121 pts) \\
  $\DK$ & $\alpha$ + $\sigma$ auto & Produit de noyaux & 1D (11 pts) \\
  \bottomrule
\end{tabular}
\caption{Comparaison des deux stratégies de combinaison.}
\end{table}

% ────────────────────────────────────────────────────────────────────────
\subsection{Notes méthodologiques : questions et arbitrages}

Plusieurs choix méthodologiques ont été guidés par des questions
pratiques :

\begin{itemize}[nosep]
  \item \textbf{Faut-il faire une ACP sur $Z$ ?} Pour l'instant, non :
    $Z$ est de petite dimension et doit rester directement
    interprétable. Une ACP classique sur $Z$ pourrait devenir utile dans
    des contextes à très haute dimension, mais ce cas ne se pose pas
    ici.
  \item \textbf{Pourquoi k-means pour A et PAM pour B/C ?} La stratégie
    A vit dans un espace euclidien réduit $\R^{K+p}$, où k-means est
    naturel. Les stratégies B et C travaillent uniquement avec des
    matrices de distances (issues de combinaisons ou de noyaux) :
    PAM (k-médoïdes) est alors adapté car il ne requiert pas de
    coordonnées explicites et accepte des distances non euclidiennes.
  \item \textbf{Comment assurer l'équité entre sources dans B/C ?}
    Côté fonctionnel, $D_0$ et $D_1$ sont normalisées dans $[0,1]$ puis
    combinées en $\Dp(\alpha)$; côté vectoriel, $Z$ est standardisé
    avant calcul de $D_s$, puis éventuellement renormalisé. Les
    paramètres $\omega$ (dans $\Dw$) et les échelles $\sigma_f$,
    $\sigma_s$ (dans $\DK$) jouent ensuite le rôle de curseurs explicites
    entre fonctionnel et vectoriel.
  \item \textbf{Pourquoi pas d'ACP dans B/C ?} Les distances de B et C
    sont construites directement à partir des courbes lissées, de leurs
    dérivées et de $Z$ standardisé. L'objectif est ici de comparer des
    \emph{architectures de distances} (et des critères de sélection),
    sans introduire de couche supplémentaire d'ACP hybride.
\end{itemize}

% ============================================================================
\newpage
\section{Résultats sur quatre jeux de données}
% ============================================================================

\subsection{Protocole expérimental}

Le pipeline a été testé sur quatre jeux de données labellisés, couvrant des
domaines variés :

\begin{table}[H]
\centering
\renewcommand{\arraystretch}{1.2}
\begin{tabular}{llcccc}
  \toprule
  \textbf{Dataset} & \textbf{Courbes $X(t)$} & $n$ & \textbf{$Z$} & $k$ & \textbf{Classes} \\
  \midrule
  Canadian Weather & Température (365j) & 35 & lat, lon, précip & 4 & Régions clim. \\
  Growth & Taille enfant (31 âges) & 93 & taille fin., croiss. & 2 & Garçon / Fille \\
  AEMET & Température Espagne (365j) & 73 & alt, lat, lon & 4 & Zones clim. \\
  Tecator & Spectres NIR (100$\lambda$) & 215 & water, protein & 3 & Niv.\ de gras \\
  \bottomrule
\end{tabular}
\caption{Les quatre jeux de données utilisés pour la validation.}
\end{table}

Pour chaque dataset, on applique le pipeline complet : lissage B-splines +
GCV, FPCA, calcul de toutes les distances, grid search sur $(\alpha, \omega)$
avec PAM, évaluation par silhouette (intrinsèque) et ARI (extrinsèque,
contre la vérité terrain).

\subsection{Le tableau principal}

\begin{table}[H]
\centering
\resizebox{\textwidth}{!}{%
\renewcommand{\arraystretch}{1.3}
\begin{tabular}{l l cc cc cc cc cc}
  \toprule
  & & \multicolumn{2}{c}{\textbf{Baseline $D_0$}} & \multicolumn{2}{c}{\textbf{Baseline $D_s$}}
    & \multicolumn{2}{c}{\textbf{A (FPCA+Z)}} & \multicolumn{2}{c}{\textbf{B ($\Dw$)}} & \multicolumn{2}{c}{\textbf{C ($\DK$)}} \\
  \cmidrule(lr){3-4}\cmidrule(lr){5-6}\cmidrule(lr){7-8}\cmidrule(lr){9-10}\cmidrule(lr){11-12}
  \textbf{Dataset} & $n$ / $k$ & Sil. & ARI & Sil. & ARI & Sil. & ARI & Sil. & ARI & Sil. & ARI \\
  \midrule
  Canadian Weather & 35 / 4 & 0.285 & 0.229 & 0.448 & 0.616 & 0.395 & \textbf{0.748} & 0.448 & 0.616 & 0.329 & 0.686 \\
  Growth           & 93 / 2 & 0.402 & 0.115 & 0.554 & 0.424 & 0.360 & \textbf{0.682} & 0.554 & 0.424 & 0.261 & 0.368 \\
  AEMET            & 73 / 4 & 0.464 & 0.167 & 0.423 & 0.325 & \textbf{0.445} & \textbf{0.367} & 0.464 & 0.167 & 0.365 & 0.316 \\
  Tecator          & 215 / 3 & 0.509 & 0.151 & 0.476 & 0.546 & 0.403 & \textbf{0.439} & 0.509 & 0.151 & 0.266 & 0.371 \\
  \bottomrule
\end{tabular}%
}

\caption{Résultats complets sur les 4 datasets (seed = 42).
Pour la stratégie B, le couple $(\alpha, \omega)$ est sélectionné par silhouette (critère non supervisé).
Les valeurs \textbf{en gras} indiquent le meilleur ARI parmi les 5 approches pour chaque dataset.
La stratégie A obtient le meilleur ARI sur 3 datasets sur 4 ; la silhouette de B retombe
systématiquement sur les baselines ($\omega \in \{0, 1\}$).}
\label{tab:resultats}
\end{table}

\subsection{Observations majeures}

\paragraph{1. $D_0$ seul est toujours insuffisant.}
L'ARI de $D_0$ varie entre 0.115 (Growth) et 0.464 (AEMET), mais est
systématiquement inférieur à au moins une autre approche sur chaque dataset.
Le niveau des courbes seul ne discrimine pas les classes.

\paragraph{2. La stratégie A est la meilleure sur 3 datasets sur 4.}
La concaténation FPCA + $Z$ avec k-means obtient le meilleur ARI sur
Canadian (0.748), Growth (0.682) et Tecator (0.439). Sur AEMET, la stratégie
A est également la meilleure approche parmi les cinq (ARI = 0.367 contre
0.325 pour $D_s$ seul).

\paragraph{3. $\alpha$ optimal est toujours $\geq 0.6$ (ARI-optimal).}

\begin{figure}[H]
\centering
\begin{tikzpicture}
  \draw[thick, ->] (0,0) -- (11, 0) node[right] {$\alpha$};
  \foreach \x in {0,1,...,10} {
    \draw (\x, -0.1) -- (\x, 0.1);
    \node[below=4pt, font=\tiny] at (\x, -0.1) {\pgfmathparse{\x/10}\pgfmathprintnumber[fixed, precision=1]{\pgfmathresult}};
  }
  % Marques
  \fill[func] (6, 0.4) circle (4pt) node[above=4pt, font=\footnotesize] {CAN};
  \fill[func] (9, 0.4) circle (4pt) node[above=4pt, font=\footnotesize] {AEM};
  \fill[func] (9, -0.5) circle (4pt) node[below=4pt, font=\footnotesize] {GRO};
  \fill[func] (10, 0.4) circle (4pt) node[above=4pt, font=\footnotesize] {TEC};

  % Zone
  \draw[very thick, mix!60, decorate, decoration={brace, amplitude=8pt, mirror}]
    (5.5, -1) -- (10.5, -1) node[midway, below=10pt, font=\small\bfseries, mix!70!black]
    {Zone optimale};
\end{tikzpicture}
\caption{Les $\alpha$ ARI-optimaux tombent tous dans $[0.6, 1.0]$.
Les dérivées ($D_1$) sont \textbf{systématiquement informatives}.}
\end{figure}

\paragraph{4. La silhouette de B retombe sur les baselines sur 4/4 datasets.}
La silhouette sélectionne $\omega \in \{0, 1\}$ sur les quatre datasets :
soit $D_s$ seul (Canadian, Growth), soit $D_0$ seul (AEMET, Tecator avec
$\omega = 1$). L'architecture $\Dw$ fonctionne — c'est le \textbf{critère de
sélection} $\alpha, \omega$ qui est le verrou.

% ============================================================================
\newpage
\section{Le paradoxe de la silhouette}
% ============================================================================

\subsection{Ce que la silhouette mesure}

La silhouette évalue si chaque point est plus proche de son cluster que du
cluster voisin. C'est un critère de \textbf{compacité géométrique} : il
récompense les clusters ronds et bien séparés.

\subsection{Pourquoi elle se trompe sur $(\alpha, \omega)$}

\begin{table}[H]
\centering
\renewcommand{\arraystretch}{1.3}
\begin{tabular}{p{5.5cm} p{5.5cm}}
  \toprule
  \textbf{Sil-optimal : $(\alpha{=}0, \omega{=}0)$} &
  \textbf{ARI-optimal : $(\alpha{=}0.6, \omega{=}0.8)$} \\
  \midrule
  Distance $= \tilde{D}_s$ (que du vectoriel) &
  Distance $= 80\%$ fonctionnel + $20\%$ vectoriel \\[4pt]
  Clusters par \textbf{géographie} &
  Clusters par \textbf{dynamique climatique} \\[4pt]
  Géométriquement \textbf{compacts} &
  Géométriquement \textbf{allongés} \\[4pt]
  Silhouette $= 0.448$ \textcolor{mix}{$\boldsymbol{\checkmark}$} &
  Silhouette $= 0.312$ \\[4pt]
  ARI $= 0.616$ &
  ARI $= 0.898$ \textcolor{mix}{$\boldsymbol{\checkmark}$} \\[2pt]
  \small 9 stations mal classées &
  \small \textbf{1 seule} station mal classée \\
  \bottomrule
\end{tabular}
\caption{Canadian Weather : la silhouette préfère des clusters \emph{ronds}
à des clusters \emph{justes}. L'ARI $= 0.898$ est obtenu en maximisant l'ARI
sur la grille $(\alpha, \omega)$ — et non par sélection silhouette. Le couple
ARI-optimal classe correctement 34 stations sur 35.}
\end{table}

\begin{keybox}[Le paradoxe résumé]
La silhouette préfère des clusters \textbf{géométriquement compacts} à des
clusters \textbf{sémantiquement corrects}. Elle est un bon critère pour
choisir $k$ (nombre de clusters), mais un mauvais critère pour choisir les
\textbf{paramètres de distance} $(\alpha, \omega)$.
\end{keybox}

\subsection{Conséquence}

\begin{resultbox}[Conclusion méthodologique]
Le problème n'est \textbf{pas} l'architecture des distances — elle atteint
ARI $\approx 0.9$ avec les bons paramètres. Le verrou principal est le
\textbf{critère de sélection} des paramètres $(\alpha, \omega)$ en contexte
purement non supervisé. La section~\ref{sec:perspectives} détaille les
pistes de recherche envisagées.
\end{resultbox}

% ============================================================================
\newpage
\section{Apports par rapport aux travaux précédents}
% ============================================================================

\begin{table}[H]
\centering
\renewcommand{\arraystretch}{1.3}
\begin{tabular}{p{5.8cm} p{5.8cm}}
  \toprule
  \textbf{Hugo Nguyen (2023)} & \textbf{Travaux actuels (2026)} \\
  \midrule
  $D_0 + D_1 \to D_p$ \newline
  (2 distances fonctionnelles) &
  $D_0 + D_1 \to D_p(\alpha)$ \newline
  (idem, \textbf{bug $D_1$ corrigé}) \\[6pt]

  \textcolor{vect}{\ding{55}} Variables $Z$ jamais utilisées &
  \textcolor{mix}{\ding{51}} $Z$ intégré dans $\Dw$ et $\DK$ \\[4pt]

  \textcolor{vect}{\ding{55}} $\Dw$ jamais implémenté &
  \textcolor{mix}{\ding{51}} $\Dw(\alpha, \omega)$ créé + grid search 2D \\[4pt]

  \textcolor{vect}{\ding{55}} $\DK$ jamais implémenté &
  \textcolor{mix}{\ding{51}} $\DK(\alpha)$ créé (noyaux gaussiens) \\[4pt]

  k-means sur matrice de distance \newline
  (traite les lignes comme des coordonnées) &
  PAM (correct pour le clustering \newline
  à partir de matrices de distance) \\[4pt]

  1 dataset (SHOM, non labellisé) &
  4 datasets labellisés + pipeline générique \\
  \bottomrule
\end{tabular}
\caption{Comparaison des contributions.}
\end{table}

\begin{resultbox}[Découverte principale]
Les dérivées ($D_1$) sont \textbf{systématiquement informatives} :
$\alpha_{\text{optimal}} \in [0.6,\, 1.0]$ sur les~4 datasets. Hugo ne
pouvait pas observer ce pattern car il n'avait pas mis en place de
recherche systématique sur $\alpha$.
\end{resultbox}

% ============================================================================
\newpage
\section{Perspectives de recherche}
\label{sec:perspectives}
% ============================================================================

L'architecture des distances a été validée sur quatre jeux de données
labellisés. Le verrou principal identifié est la \textbf{sélection
automatique des paramètres} $(\alpha, \omega)$ en contexte purement non
supervisé, où l'ARI n'est pas disponible. En non supervisé pur, il faut
résoudre \textbf{deux problèmes couplés simultanément} :
\begin{enumerate}[nosep]
  \item Choisir $k$ (nombre de clusters).
  \item Choisir $(\alpha, \omega)$ (paramètres de distance).
\end{enumerate}
Ces problèmes sont couplés car le meilleur $(\alpha, \omega)$ dépend de $k$,
et inversement. Voici les pistes de recherche envisagées.

% ────────────────────────────────────────────────────────────────────────
\subsection{Sélection par stabilité du clustering}

La piste la plus prometteuse repose sur une idée simple : le bon triplet
$(k, \alpha, \omega)$ est celui qui produit des clusters
\textbf{reproductibles} lorsqu'on perturbe les données.

\paragraph{Protocole.}
Pour chaque triplet $(k, \alpha, \omega)$ dans la grille :
\begin{enumerate}[nosep]
  \item Calculer la partition $\mathcal{C}$ sur les $n$ individus complets
    (PAM avec la distance $\Dw(\alpha, \omega)$).
  \item Répéter $B$ fois ($B \approx 50$) :
    \begin{enumerate}[nosep]
      \item Sous-échantillonner $80\%$ des individus aléatoirement.
      \item Recalculer la distance et la partition $\mathcal{C}_b^*$ sur ce
        sous-échantillon.
      \item Mesurer l'ARI entre $\mathcal{C}$ (restreinte aux individus
        présents) et $\mathcal{C}_b^*$.
    \end{enumerate}
  \item La \textbf{stabilité} du triplet =
    $\text{Stab}(k, \alpha, \omega) = \frac{1}{B}\sum_{b=1}^{B}
    \text{ARI}(\mathcal{C}, \mathcal{C}_b^*)$.
\end{enumerate}
Le triplet $(k^*, \alpha^*, \omega^*)$ avec la stabilité maximale est retenu.

\paragraph{Pourquoi la stabilité ?}
\begin{itemize}[nosep]
  \item Pas de biais vers des formes géométriques particulières
    (contrairement à la silhouette).
  \item Fonctionne aussi pour sélectionner $k$ : un $k$ trop grand ou trop
    petit produit des clusters instables.
  \item Ne nécessite \textbf{aucun label} — l'ARI est calculé entre deux
    partitions produites par la méthode elle-même.
  \item Mesure une propriété \textbf{intrinsèquement désirable} : la
    reproductibilité.
\end{itemize}

\paragraph{Coût calculatoire.}
Pour une grille $(\alpha, \omega) \in \{0, 0.1, \ldots, 1\}^2$,
$k \in \{2, \ldots, 6\}$, et $B = 50$ : cela représente
$11 \times 11 \times 5 \times 50 = 30\,250$ exécutions de PAM. C'est
un coût important mais tout à fait faisable (PAM sur $n < 300$ est rapide).

% ────────────────────────────────────────────────────────────────────────
\subsection{Vote entre critères de validation}

Plutôt que de s'appuyer sur un unique critère (silhouette, gap, ou
stabilité), on peut faire \textbf{voter} plusieurs critères pour élire le
meilleur $(k, \alpha, \omega)$.

\paragraph{Principe.}
Pour chaque triplet $(k, \alpha, \omega)$, on calcule $M$ critères
de validation interne :
\begin{itemize}[nosep]
  \item \textbf{Silhouette} : compacité et séparation géométrique.
  \item \textbf{Gap statistic} \cite{tibshirani2001} : écart d'inertie
    par rapport à une référence uniforme.
  \item \textbf{Prediction strength} \cite{tibshirani2005} :
    reproductibilité par validation croisée.
  \item \textbf{Stabilité} (bootstrap, cf.\ section précédente).
  \item \textbf{Indice de Calinski-Harabasz} \cite{calinski1974} : rapport variance
    inter/intra.
  \item \textbf{Indice de Davies-Bouldin} \cite{davies1979} : ratio de dispersion intra par
    rapport à la séparation inter.
\end{itemize}

Chaque critère $m$ produit un \textbf{classement} des triplets
$(k, \alpha, \omega)$. On agrège ces classements par une règle de vote :

\begin{itemize}[nosep]
  \item \textbf{Vote majoritaire} : le triplet qui arrive en tête du plus
    grand nombre de critères.
  \item \textbf{Score de Borda} : chaque critère attribue un rang, on somme
    les rangs, le triplet avec la somme de rangs la plus faible gagne.
  \item \textbf{Vote pondéré} : pondérer les critères par leur fiabilité
    estimée (par exemple sur les datasets labellisés de validation).
\end{itemize}

L'avantage du vote est qu'il réduit la \textbf{dépendance à un seul
indicateur}. Un critère qui se trompe (comme la silhouette sur $\omega$)
sera corrigé par les autres, à condition que la majorité des critères
capturent un signal pertinent.

% ────────────────────────────────────────────────────────────────────────
\subsection{Optimisation bayésienne}

Le grid search 3D $(k, \alpha, \omega)$ évalue systématiquement toutes les
combinaisons, ce qui est coûteux. L'\textbf{optimisation bayésienne}
propose une alternative intelligente :

\begin{enumerate}[nosep]
  \item Évaluer un petit nombre de triplets initiaux (design latin
    hypercube, $\sim 15$ points).
  \item Ajuster un \textbf{processus gaussien} (GP) sur la surface
    $S(k, \alpha, \omega)$ (critère de stabilité ou score de vote).
  \item Choisir le prochain triplet à évaluer par un critère
    d'\textbf{acquisition} (Expected Improvement) : explorer là où le GP
    prédit un gain potentiel ou une forte incertitude.
  \item Itérer jusqu'à convergence ($\sim 20$--$30$ évaluations au total).
\end{enumerate}

Cela réduit le nombre d'évaluations de $\sim 600$ (grid search complet) à
$\sim 30$, rendant l'approche applicable même à de grands jeux de données.

% ────────────────────────────────────────────────────────────────────────
\subsection{Méta-calibration sur données labellisées}

Les résultats montrent que $\alpha_{\text{optimal}} \geq 0.6$ sur les~4
datasets testés (cf.~Table~\ref{tab:resultats}). Ce pattern pourrait être
exploité pour \textbf{fixer} $\alpha$ à une valeur méta-calibrée, réduisant
le problème 3D $(k, \alpha, \omega)$ à un problème 2D $(k, \omega)$. La
robustesse de cette approche nécessite d'être validée sur un plus grand
nombre de jeux de données fonctionnels labellisés (idéalement 15--20).

% ────────────────────────────────────────────────────────────────────────
\subsection{Enrichissement de l'architecture des distances}

\begin{itemize}
  \item \textbf{Dérivées d'ordre supérieur} : $D_2$ (accélération) pourrait
    capturer des changements de courbure informatifs sur certains types de
    données.
  \item \textbf{Combinaison de $\Dw$ et $\DK$} : $\Dw$ (somme pondérée) et
    $\DK$ (produit de noyaux) produisent des partitions différentes.
    Explorer des stratégies d'ensemble qui agrègent les partitions issues de
    plusieurs distances (\textbf{consensus clustering}).
\end{itemize}

% ────────────────────────────────────────────────────────────────────────
\subsection{Vers une HFV-PCA (hybride fonctionnelle + vectorielle)}

L'approche développée dans ce rapport repose sur une \textbf{RS-PCA} :
les fonctions sont d'abord résumées par FPCA, puis concaténées avec les
variables vectorielles $Z$ pour travailler dans un espace réduit
$\R^{K+p}$, sur lequel on définit des distances adaptées au clustering.
Une extension naturelle consiste à formuler l'analyse dans un espace de
Hilbert \textbf{hybride} $\mathcal{H}_f \times \R^p$, où chaque
composante principale est elle-même un objet mixte (fonction +
vecteur). C'est l'idée des approches de type \textbf{HFV-PCA}
(Hybrid Functional and Vector PCA).

Dans ce cadre, les composantes principales seraient définies comme les
vecteurs propres d'un opérateur de covariance sur
$\mathcal{H}_f \times \R^p$, produisant un jeu unique de scores
hybrides pour chaque individu. Ce chantier méthodologique dépasse le
cadre de ce rapport, centré sur l'architecture de distances et les
critères de sélection pour le clustering, mais il constitue une piste
prioritaire pour des travaux ultérieurs.

% ────────────────────────────────────────────────────────────────────────
\subsection{Application aux données non labellisées}

Le test final de l'architecture sera son application à des données
réelles non labellisées (par exemple, les profils de célérité du son du
SHOM). L'enjeu est triple : (1)~déterminer automatiquement le nombre de
clusters~$k$, (2)~sélectionner $(\alpha, \omega)$ sans vérité terrain, et
(3)~valider la pertinence des clusters obtenus par expertise métier.
Les pistes de recherche présentées ci-dessus (stabilité, vote entre
critères, optimisation bayésienne, méta-calibration) visent précisément
à résoudre ce verrou.

% ============================================================================
\begin{thebibliography}{9}

\bibitem{tibshirani2001}
  R.~Tibshirani, G.~Walther et T.~Hastie,
  \textit{Estimating the number of clusters in a data set via the gap statistic},
  Journal of the Royal Statistical Society B, 63(2):411--423, 2001.

\bibitem{tibshirani2005}
  R.~Tibshirani et G.~Walther,
  \textit{Cluster validation by prediction strength},
  Journal of Computational and Graphical Statistics, 14(3):511--528, 2005.

\bibitem{calinski1974}
  T.~Calinski et J.~Harabasz,
  \textit{A dendrite method for cluster analysis},
  Communications in Statistics, 3(1):1--27, 1974.

\bibitem{davies1979}
  D.L.~Davies et D.W.~Bouldin,
  \textit{A cluster separation measure},
  IEEE Transactions on Pattern Analysis and Machine Intelligence, 1(2):224--227, 1979.

\end{thebibliography}

\end{document}
